% Generated by support/export-artifacts.py; do not edit.
% Slug: ch76-errors
% Source: chapters/ch76.tex:132
% TikZ source: chapters/ch76.tex:134
% Original label: \label{fig:ch76-errors}
% Original caption (LaTeX source):
%   A common native status maps back to each runtime's own error mechanism.
% Caption: A common native status maps back to each runtime's own error mechanism.
% Setup: support/compiler-pseudocode.sty (TikZ libraries and shared styles).
\begin{tikzpicture}
\node[basic block] (native) at (0,0) {native visitor returns status};
\node[basic block,text width=4cm] (python) at (-3,-1.6) {Python: null plus exception};
\node[basic block,text width=4cm] (lua) at (3,-1.6) {Lua: re-raise retained object};
\draw[cfg edge] (native)--(python); \draw[cfg edge] (native)--(lua);
\end{tikzpicture}

% Original figure LaTeX:
% \begin{figure}[htbp]
% \centering
% \begin{tikzpicture}
% \node[basic block] (native) at (0,0) {native visitor returns status};
% \node[basic block,text width=4cm] (python) at (-3,-1.6) {Python: null plus exception};
% \node[basic block,text width=4cm] (lua) at (3,-1.6) {Lua: re-raise retained object};
% \draw[cfg edge] (native)--(python); \draw[cfg edge] (native)--(lua);
% \end{tikzpicture}
% \caption{A common native status maps back to each runtime's own error mechanism.}
% \label{fig:ch76-errors}
% \end{figure}
